\chapter*{Preface\sv{Förord}}
\begin{instructions}
Tell the reader about any collaboration and about the contribution of each of the authors.

Tell the reader whether you want to thank someone or acknowledge something. For example, if you have used any help (automatic or from a person) to improve your text or your work, this should be acknowledged here.

This is the only place where you can mention the project owner by name (company name/teacher/entity).
\end{instructions}

\section*{Use of Generative AI\sv{Användning av generativ AI} \hint{(mandatory)}}
\label{sec:genai}
\begin{instructions}
This section is \textbf{mandatory}. You must describe how generative AI (e.g.\ ChatGPT, Copilot, Claude, Gemini or similar tools) has been used during the thesis project. Describe
\begin{itemize}
    \item which tools you have used (name and, if known, version),
    \item what you used them for, for example programming, debugging, searching for or summarising literature, generating ideas, translating, or improving the language of the text,
    \item in which parts of the work and of the report they were used,
    \item how you checked and took responsibility for the output (you are responsible for everything in the thesis, including errors introduced by AI tools).
\end{itemize}
If you have not used generative AI at all, state that explicitly in this section.
\end{instructions}
